%% generated by tools/make_tables.py -- do not edit
\begin{table}[t]
\centering
\caption{光测地线积分器的全部数值判据。左列为源码中的字面写法，右列为物理含义；本表由 \texttt{tools/make\_tables.py} 直接匹配 \texttt{web/src/shaders.js} 的正则生成，任一常量被改动而未被本表捕获时编译脚本会立即报错。}
\label{tab:integrator}
\scriptsize
\begin{tabular}{p{0.34\linewidth}p{0.58\linewidth}}
\toprule
源码字面量 & 作用与取值 \\
\midrule
\texttt{for (int i = 0; i < N; i++)} & 编译期迭代硬上限（防止病态光线挂死 GPU） \\
\texttt{uMaxSteps} & 运行期步数上限 $n_{\max}$，用户可在 60--900 间调节 \\
\texttt{u > 0.5} & 捕获判据：$r<\rs$ 后不再有可能逃逸 \\
\texttt{u < uMin \&\& du < 0} & 逃逸判据：向外运动且已越过 $R_{\rm esc}$ \\
\texttt{uMin = 1/max(uEscapeR, 1.6 * r0)} & 逃逸半径至少为观测者半径的 1.6 倍，避免相机贴视界时的伪逃逸 \\
\texttt{hmax} & 步长上界：远场 $0.35$ rad，近视界收窄到 $0.05$ rad \\
\texttt{h = uTol * u / |du|} & 把 $|\Delta u/u|$ 限制在 $tol$ 量级（防止除零） \\
\texttt{clamp(h, 0.002, hmax)} & 步长下界固定为 $0.002$ rad：转折点附近 $|\dd u|\to0$ 时不无限细分 \\
\texttt{uStepScale} & 全局步长缩放 $\sigma_h$，与 $1/tol$ 近似等效（见 \S\ref{sec:perf}） \\
\texttt{btan < 1e-5} & 径向退化光线（$\sin\psi\to0$）走解析分支，不进入积分循环 \\
\texttt{int K = nearDisk ? 5 : 1} & 近盘段的一维子步数，用于把体积发射积到亚像素精度 \\
\texttt{dtau = min(\ldots, 6.0)} & 单段光学厚度上限，防止 $\exp(-\tau)$ 下溢导致透明度突跳 \\
\texttt{trans > 0.002} & 透射率截断：不透明后不再采样盘，节省算力 \\
\texttt{u < 1e-6} & 数值上「跑到无穷远」的兜底退出 \\
\bottomrule
\end{tabular}

\end{table}
